\documentclass[10pt, a4paper, oneside]{ctexart}
\usepackage{amsmath, amsthm, amssymb, bm, xcolor, framed, graphicx, hyperref, mathrsfs,ulem,float}
\usepackage{titling,geometry}
\usepackage{graphicx,caption,subcaption,enumitem,wrapfig}
\usepackage{lipsum,zhlipsum}
\usepackage[siunitx]{circuitikz}
\geometry{a4paper,scale=0.75}
\setlength{\droptitle}{-1.5cm}

\colorlet{shadecolor1}{blue!5} %颜色可自定义
\colorlet{framecolor1}{shadecolor1}

\colorlet{shadecolor2}{shadecolor1!50!black!10}
\colorlet{framecolor2}{shadecolor2}

\newenvironment{framedshaded}[2]
  {\def\FrameCommand{\setlength{\FrameRule}{1pt}\fboxsep=\FrameSep \fcolorbox{#1}{#2}} \MakeFramed {\FrameRestore}}
  {\endMakeFramed}

\newcounter{problemname}
\newcounter{subproblemname}

\newenvironment{problem} [1][\arabic{problemname}]
  {\begin{framedshaded}{framecolor1}{shadecolor1}\stepcounter{problemname}\par\noindent\textbf{#1. }}
  {\end{framedshaded}\par\setcounter{subproblemname}{0}}

\newenvironment{problem*}[1][\arabic{problemname}]
  {\begin{framedshaded}{framecolor1}{shadecolor1}\par\noindent\textbf{#1. }}
  {\end{framedshaded}\par\setcounter{subproblemname}{0}}

\newenvironment{subproblem}[1][\arabic{problemname}.\arabic{subproblemname}]
  {\begin{framedshaded}{framecolor2}{shadecolor2}\stepcounter{subproblemname}\par\noindent\textbf{#1. }}
  {\end{framedshaded}\par}

\newenvironment{subproblem*}[1][\arabic{problemname}.\arabic{subproblemname}]
  {\begin{framedshaded}{framecolor2}{shadecolor2}\par\noindent\textbf{#1. }}
  {\end{framedshaded}\par}

\newenvironment{solution}
  {\par\noindent\textbf{解 }}
  {\par }

\renewenvironment{proof}
  {\par\noindent\textbf{证 }}
  {\hfill $\blacksquare$ \par }

\newenvironment{note}[1][题目\arabic{problemname}的注记]
  {\par\noindent\textbf{#1. }}
  {\par}

\renewcommand{\figurename}{fig}
\renewcommand{\thefigure}{\arabic{figure}}
\captionsetup{margin=10pt , font=small, labelfont=bf, labelsep=period,skip=5pt , labelformat=simple}


\title{\textbf{作业模板}}
\author{作者\textbackslash 解答者}
\date{}

\begin{document}
\maketitle
\vspace*{-1.0cm}

这是一个LaTeX作业排版模板的使用介绍。模板提供了方便的环境来排版作业题目、子题目、解答和证明等。

模板主要环境包括：
\begin{itemize}
  \item \verb|problem|：用于主题目，会自动编号。
  \item \verb|subproblem|：用于子题目，在主题目下自动编号。
  \item \verb|solution|：用于书写解答。
  \item \verb|proof|：用于书写证明，结尾自动添加证明结束符号。
  \item \verb|note|：用于添加注记，与当前题目关联。
\end{itemize}
使用自定义编号时，可以在环境开始命令中指定，例如 \verb|\begin{problem}[自定义编号]|。使用带星号的环境（如 problem*）可以避免编号自增。

模板还预设了颜色和框架样式，使内容更美观。支持插入图片和 TikZ 绘图，示例中展示了基本用法。

要使用模板，将提供的代码保存为 .tex 文件，使用支持中文的 LaTeX 引擎（如 XeLaTeX）编译。然后根据需要修改内容，替换示例文字为自己的题目和解答。

参考模板中的示例来编写题目、子题目、解答和注记。对于更复杂的排版，可以调整模板中的宏包和设置。

\begin{problem}
  这是一个好的题目，因为它使用了\verb|problem|环境。
\end{problem}
\begin{subproblem}
  因为这是一个好题目，所有它有一道子题目，且使用了\verb|subproblem|环境。
\end{subproblem}
\begin{solution}
  为了表达解决这个好问题的喜悦，我使用\texttt{Ti\textit{k}Z}绘制一张好图(fig.\ref{fig:nicefig})。
  \begin{note}
    Ti\textit{k}Z是``Ti\textit{k}Z ist \textit{k}ein Zeichenprogramm''的递归缩写。
  \end{note}
\end{solution}
\begin{figure}[H]
  \centering
  \begin{tikzpicture}
    \filldraw[rounded corners,fill=blue!10] (0,0) rectangle (4,2);
    \draw (1.2,1) -- (2,0.2) -- (3.5,1.8);
    \node[shape=rectangle] (txt) at (2, 1) {\small{Hello!}};
  \end{tikzpicture}
  \caption{好图（也可以称之为``不错的图片''）。}\label{fig:nicefig}
\end{figure}
\begin{subproblem*}[题目\arabic{problemname}的附加子题]
  使用\verb|\begin{problem}[自定义题目编号]|可以自定义题目编号。默认后续题目会继续递增题目编号记数器，使用\verb|problem*|可以使编号不自增。对于\verb|subproblem|也可以如此使用。
\end{subproblem*}
\begin{subproblem}
  \footnotesize{一段乱数假文。\lipsum[1]}
\end{subproblem}
\begin{solution}
  \footnotesize{也是乱数假文。\zhlipsum[1]}
\end{solution}

\begin{problem}[2.正经例题]
	\noindent 设计一个 \( R = 2\,\text{k}\Omega \)，\( V_{OL} = 0.05\,\text{V} \) 的电阻负载反相器，增强型 nMOS 驱动晶体管参数如下：
	\begin{flalign*}
	V_{DD} &= 1.1\,\text{V} &\\
	V_{T0} &= 0.52\,\text{V} &\\
	\gamma &= 0\,\text{V}^{1/2}& \\
	\lambda &= 0 &\\
	\mu_n C_{ox} &= 216\,\mu\text{A/V}^2 &
	\end{flalign*}
	\noindent \textbf{a.} 求所需的宽长比 \( W/L \)。

	\noindent \textbf{b.} 求 \( V_{IL} \) 和 \( V_{IH} \)。

\noindent \textbf{c.} 求噪声容限 \( NM_L \) 和 \( NM_H \)。
\end{problem}
\begin{solution}
\begin{itemize}
	\item[\textbf{a.\ }]输入为 \( V_{DD} \) 时，输出为 \( V_{OL} \)，可列写：
		\[
		\frac{V_{DD} - V_{OL}}{R} = \frac{1}{2}\mu_n C_{ox} \frac{W}{L} \left[ 2(V_{DD} - V_{T0})V_{OL} - V_{OL}^2 \right]
		\]
		解得 \( {W}/{L} = 87.59 \)
	\item[\textbf{b.\ }] \( k_n R_L = 37.84 \, V^{-1} \)
		\[
		V_{IL} = V_{T0} + \frac{1}{k_n R_L} = 0.546 \, \text{V}
		\]
		\[
		V_{IH} = V_{T0} + \sqrt{\frac{8}{3} \cdot \frac{V_{DD}}{k_n R_L}} - \frac{1}{k_n R_L} = 0.772 \, \text{V}
		\]
	\item[\textbf{c.\ }] \[ NM_L = V_{IL} - V_{OL} = 0.496 \, \text{V} \]
		\[ NM_H = V_{OH} - V_{IH} = 0.328 \, \text{V} \]
\end{itemize}
\end{solution}
\end{document}